\documentclass[11pt]{scrartcl}
\usepackage[headheight=1pt, includeheadfoot, headsep=1cm, top=1.5cm, bottom=1.3cm, left=2cm, right=2cm]{geometry}
\usepackage{mathtools}
\usepackage{amsmath,amsthm, amscd, amssymb, amsfonts}
\usepackage[english]{babel}
\usepackage{fancyhdr}
\usepackage{enumitem}
\setlist{itemsep=0.5em}
\usepackage{xcolor}
\definecolor{forestgreen}{rgb}{0.13, 0.55, 0.13}
\usepackage[colorlinks = true,
linkcolor = forestgreen,
urlcolor  = forestgreen,
citecolor = forestgreen,
anchorcolor = forestgreen]{hyperref}
\usepackage{multicol}
\usepackage{multienum}
\usepackage{euler}
\usepackage[OT1]{eulervm}
\renewcommand{\rmdefault}{pplx}
\usepackage{tikz}
\usepackage{mathabx}
%\usepackage{lastpage}





%% header, footer definitions
\def\course{MATH 314}
\def\courselink{https://www.gsanmarco.com/graph-theory}
\def\assignment{HW \#1}
\def\assignlink{hw1.tex}
\def\due{Due on Jan 24}




\pagestyle{fancy}
\fancyhead[L]{\course}
\fancyhead[C]{\assignment} 
\fancyhead[R]{\due} 
%\renewcommand{\footrulewidth}{0.4pt}
%\fancyfoot{}
%\fancyfoot[R]{\fontsize{8}{10} \selectfont  Page \thepage~of \pageref{LastPage}}
%\fancyfoot[L]{\fontsize{10}{12} \selectfont  Source file available at \href{\courselink}{\textbf{\courselink}}}

%% environments
% theorem style
\newtheorem{theorem}{Theorem}
\newtheorem{lemma}[theorem]{Lemma}
\newtheorem{corollary}[theorem]{Corollary}
\newtheorem{claim}[theorem]{Claim}
\newtheorem{defn}[theorem]{Definition}

% definition style
\theoremstyle{definition}
\newtheorem{problem}{Problem}

% custom
\newenvironment{solution}{\noindent\textbf{Solution.}}{\qed}

%% custom commands

\newcommand*{\abs}[1]{\lvert #1\rvert}     % absolute values, cardinality
\newcommand*{\eps}{\varepsilon}            % prettier epsilon
\newcommand*{\R}{\mathbb{R}}               % real numbers
\newcommand*{\C}{\mathbb{C}}               % complex numbers
\newcommand*{\Q}{\mathbb{Q}}               % rational numbers
\newcommand*{\Z}{\mathbb{Z}}               % integers
\newcommand*{\N}{\mathbb{N}}               % natural numbers
\newcommand*{\F}{\mathbb{F}}               % field
\newcommand*{\family}{\mathcal{F}}         % family
\newcommand*{\GL}{\mathbf{GL}}             % general linear group
\newcommand*{\what}[1]{\widehat{#1}}       % wider hat
\newcommand*{\wbar}[1]{\overline{#1}}      % wider bar
\newcommand*{\mcal}[1]{\mathcal{#1}}       % mathcal
\newcommand*{\mbf}[1]{\mathbf{#1}}         % mathbf
\newcommand*{\mbb}[1]{\mathbb{#1}}         % mathbb

% operators
\let\Pr\relax\DeclareMathOperator*{\Pr}{\mathbf{Pr}}       % probability
\DeclareMathOperator*{\E}{\mathbb{E}}          % expectation
\DeclareMathOperator*{\Var}{\mathbf{Var}}      % variance
\DeclareMathOperator*{\Span}{span}             % span
\DeclareMathOperator{\tr}{tr}                  % trace
\DeclareMathOperator{\rank}{rank}              % rank
\DeclareMathOperator{\supp}{supp}              % support
\DeclareMathOperator{\diam}{diam}



\begin{document}
\noindent  Unless explicitly requested by a problem, do not include sketches as part of your proof.
You are free to use the result from any problem on this (or previous) assignment as a part of your solution to a different problem even if you have not solved the former problem.

\vskip10pt

\begin{problem}[2 pts]
Let $S$ be a finite set of integers and let $G$ be the graph on vertex set $S$ where for any $x,y\in S$, $xy\in E(G)$ if and only if $x+y$ is odd.
\begin{enumerate}
\item Prove that $G$ is a bipartite graph for any such $S$.
\item If $S=\{1, 2, \dots, 100\}$, what is $\abs{E(G)}$? Justify your answer.
\end{enumerate}
\end{problem}

\begin{solution}
1. We can constrict a bipartition. Indeed, notice that $xy\in E(G)$ if and only if $x$ and $y$ have different parities. Consider thus $X=\{x\in S\colon x \text { is even}\}$ and $Y=\{x\in S\colon x \text { is odd}\}$, so that $V(G)=X\sqcup Y$. Since every edge of $G$ connects a vertex in $X$ to a vertex in $Y$, both $X$ and $Y$ are independent, so $X\sqcup Y $ is a bipartition of $G$.

2. As mentioned above, $E(G)=\{xy\colon x \in X, y \in Y\}$. Thus $\vert E(G)\vert = \vert X \vert \cdot \vert Y \vert=50 \cdot 50 = 2500$. 
\end{solution}

\hrulefill
\begin{problem}[2 pts] 
Suppose that $(u=v_0,v_1,\dots,v_k=v)$ is a $u$--$v$ geodesic.
Prove that $d(u,v_i)=i$ for all $i\in\{0,\dots,k\}$.
\end{problem}
\begin{solution}
Notice first for later use that $d(u,v)=k$, because $(u=v_0,v_1,\dots,v_k=v)$.

Fix any $i\in\{0,\dots,k\}$. Since$(u=v_0,v_1,\dots,v_i)$ is a $u$--$v_i$ of length $i$, we know that $d(u,v_i)\leq i$. Suppose for the sake of contradiction that  $d(u,v_i)=\ell< i$; hence there is a $u$--$v_i$ geodesic $(u = w_0,\dots,w_\ell = v_i)$. But then $(u = w_0,\dots,w_\ell = v_i, v_{i+1},\dots,v_k= v)$ is a $u$--$v$ walk of length $\ell+(k-i)$.  Theorem 1.6 then implies that there is a $u$--$v$ path of length at most $\ell+(k-i)<k$, contradicting the previously established fact that $d(u, v) = k$.
\end{solution}

\hrulefill

\begin{problem}[2 pts]\label{abpath}
Let $G$ be a graph.
For two non-empty subsets $A,B\subseteq V(G)$, an $A$--$B$ path is a path in $G$ which connects some vertex of $A$ to some vertex of $B$.
Prove that if $P$ is a minim\textbf{al} $A$--$B$ path, then $P$ contains exactly one vertex from $A$ and contains exactly one vertex from $B$.
\end{problem}
\begin{solution}
Let $P = (v_0,\dots,v_k)$ be a minimal $A$--$B$ path with $v_0 in A$ and $v_k \in B$. In case $k=0$ (which can happen only if A and B intersect), then we are trivially done, so we may assume that $k\geq1$. Suppose for a contradiction that there is some $i \in \{1, \dots k\}$ such that also $v_i$ is in $A$ (note that $i$ could equal $k$ if $A$ and $B$ intersect). But then $(v_0,\dots,v_k)$ is also an $A$--$B$ path which is a proper sub-path of $P$, a contradiction to the minimality of P. Therefore $P$ contains exactly one vertex of $A$, namely $v_0$. A symmetric argument establishes that $P$ contains exactly one vertex (namely $v_k$) of $B$.
\end{solution}

\hrulefill

\begin{problem}[2 pts]
Let $G$ be a connected graph.
Prove that any two maxim\textbf{um} paths in $G$ must share some vertex.
\end{problem}
\begin{solution}
Let $k$ denote the length of any maximum path in $G$. Assume for a contradiction that $(v_0,\dots, v_k)$ and $(u_0,\dots, u_k)$ are two maximum paths in $G$ with no vertices in common. Consider the subsets $A=\{v_0,\dots, v_k\}$ and $B=\{u_0,\dots, u_k\}$ of $V(G)$. Since $G$ is connected, it contains an $A$--$B$ path. Consider a minimal $A$--$B$ path $P=(w_0, \dots w_\ell)$ with $w_0\in A$ and $w_\ell \in B$. Note that $\ell \geq 1$ since $A$ and $B$ are disjoint, and that problem \ref{abpath} tells us all vertices of $P$ other than $w_0$ and $w_\ell$ belong to neither $A$ nor $B$. Since $w_0\in A$ and $w_\ell\in B$, we have $w_0=v_i$ and $w_\ell=u_j$ for some $i$, $j$. Consider the paths 
\begin{align*}
P_1&=(v_0, \dots, v_{i-1}, v_i=w_0, w_1, \dots, w_{\ell-1}, w_\ell=u_j, u_{j+1}, \dots,u_k), \\
P_2&=(v_{k}, \dots, v_{i+1}, v_i=w_0, w_1, \dots, w_{\ell-1}, w_\ell=u_j, u_{j-1}, \dots, u_0).
\end{align*}
Observe that the length of $P_1$ is $i+\ell+(k-j)\geq k+1+(i-j)$, where we used the fact that $\ell\geq 1$. Similarly, the length of $P_2$ is $(k-i)+\ell+j\geq k+1+(j-i)$
Since either $i\geq j$ or $j\geq i$, we see that one of $P_1$, $P_2$ has length at least $k+1$. This contradicts the fact the length of a maximum path in $G$ is $k$.
\end{solution}

\hrulefill
\begin{problem}[2 pts]
Prove that a graph $G$ is bipartite if and only if every subgraph $H$ of $G$ has an independent set consisting of at least half of $V(H)$.
(recall that an independent set is a set of vertices which induce no edges)
\end{problem}

\begin{solution}
Assume first that $G$ is bipartite and let $V(G)=X\sqcup Y$ be a bipartition of $G$. If $H$ is any subgraph of $G$, we discussed in class that $(V(H)\cap X) \sqcup(V(H)\cap Y)$ is a bipartition of $H$. In particular both $V(H)\cap X$ and $V(H)\cap Y$ are independet sets of $H$, and since their union gives $V(H)$, at least one of them have size $\geq \frac{V(H)}{2}$. 

Conversely, assume that $G$ is not bipartite. Then by Theorem 1.12, $G$ contains an odd cycle $C$, say of length $n$. Then $C$ is a subgraph of $G$, and we showed in class that largest size of an independent set of $C$ is $\frac{n-1}{2}$.
\end{solution}


\end{document}





